% Generated draft: requires mathematical and source review.
\documentclass[11pt]{article}
\usepackage[margin=25mm]{geometry}
\usepackage{fontspec}
\setmainfont{Latin Modern Roman}
\usepackage{amsmath,amssymb,mathtools,unicode-math}
\setmathfont{Latin Modern Math}
\usepackage{enumitem,needspace,xurl,hyperref}
\setlength{\parindent}{0pt}
\setlength{\parskip}{5pt}
\setlength{\emergencystretch}{3em}
\newcommand{\N}{\mathbb{N}}
\newcommand{\Z}{\mathbb{Z}}
\newcommand{\Q}{\mathbb{Q}}
\newcommand{\R}{\mathbb{R}}
\newcommand{\C}{\mathbb{C}}
\newcommand{\F}{\mathbb{F}}
\newcommand{\A}{\mathbb{A}}
\newcommand{\PP}{\mathbb{P}}
\newcommand{\E}{\mathbb{E}}
\newcommand{\eps}{\varepsilon}
\newcommand{\dd}{\,\mathrm{d}}
\begin{document}
\section*{UnsolvedMath ID: 30004423}
\textbf{Source problem code:} catalogue entry
\subsection{Definitions and mathematical statement}
For terminology, a partition is a finite weakly decreasing sequence of nonnegative integers, with trailing zeros ignored. In the usual English diagram convention, its cells are pairs \((i,j)\) with \(1\le j\le\lambda_i\). The arm length of a cell is
\[
 a_\lambda(i,j)=\lambda_i-j,
\]
the number of cells strictly to its right in the same row.

Under a standard convention, a partition contained in a rectangle with \(r\) rows and \(c\) columns satisfies \(c\ge\lambda_1\ge\cdots\ge\lambda_r\ge0\), and its complementary partition is
\[
 \lambda^{\mathrm c}_i=c-\lambda_{r+1-i}\qquad(1\le i\le r).
\]
This describes rotating the unoccupied cells of the rectangle through \(180\) degrees. These definitions explain standard terminology; the supplied record does not specify its diagram or complement conventions, or its permitted rectangle dimensions.

\textbf{Supplied problem request.} Find a representation-theoretic interpretation of the stronger arm-length-refined multiset identity for complementary partitions in a rectangle.

\textbf{Unresolved specification.} The supplied record does not state this multiset identity or the weaker identity relative to which it is stronger. It does not identify the multisets, their entries and multiplicities, the arm-length refinement, or the quantifiers and boundary cases. It also does not specify the representations or the required meaning of an interpretation. Consequently, the request above can be preserved, but a precise mathematical target cannot be recovered from the supplied evidence. The terminology given here does not supply the missing identity.
\subsection{Short English statement}
The problem asks for a representation-theoretic explanation of an arm-length-refined multiset identity involving complementary partitions in a rectangle. The identity itself is absent from the supplied record, so its exact content and scope remain unresolved. The supplied open status is a dated curation label, not an independently verified conclusion.
\subsection{Sources}
\begin{itemize}
\item Local UnsolvedMath registry: permanent ID \(30004423\), source code \texttt{catalogue entry}, title \emph{Arm-Length Refinements for Complementary Partitions}. The supplied original, clean, and current statements all contain the same request. No individual proposer is supplied.
\item Supplied report citation: \emph{Mini-Workshop: Kronecker, Plethysm, and Sylow Branching Coefficients and their Applications to Complexity Theory} (2021), Oberwolfach Reports, pp.\ 517--568, DOI \texttt{10.4171/owr/2020/9}. Supplied source URL and literature-source link, labelled ``original OWR report'': \url{https://doi.org/10.4171/owr/2020/9}. The background describes Oberwolfach Reports as records of mathematical workshops and their open-problem discussions. The report contents were not supplied or independently consulted here.
\item Literature assessment dated exactly \texttt{2026-08-22}: the supplied curation labels the problem open and reports that no subsequent general construction resolving this specific request was identified. This is supplied dated literature triage, not independently verified research or a claim of current status. The background explicitly requests independent verification of time-sensitive status and cited claims.
\item The supplied proposed year and year tag are 2020, while the supplied report citation gives 2021. Both are retained without resolving their relationship. The registry also supplies the label \texttt{corrected\_verified} and a claim of manual semantic, typographic, and LaTeX review; that supplied claim does not resolve the missing mathematical formulation.
\end{itemize}
\end{document}
